\chapter{Analysis}

\section{Introduction}

This chapter analyses what can reasonably be concluded from the results in
Chapter 4. It considers the evidence provided by the controlled baselines,
structural ablations, generalisation audit and candidate-search experiment. It
also distinguishes results that support the proposed design from those that
reveal uncertainty or limitations.

The analysis remains within the scope of the experiments. A model score
represents compatibility learned from the processed Polyvore data. It does not
measure whether an individual user likes an outfit, whether the outfit is
objectively fashionable or whether a recommended concept will lead to a
purchase.

\section{Compatibility performance and baseline evidence}

The selected compatibility model achieved an AUC of 0.8583 on the held-out
test data. This means that the model usually assigned a higher score to a
randomly selected positive outfit than to a randomly selected negative outfit
\parencite{fawcett2006roc}. It does not mean that 85.83\% of recommendations are
correct or will be liked by users.

The non-learned SigLIP similarity baseline reached an AUC of 0.7012. Frozen
visual features therefore contained some information related to the dataset
labels without compatibility training. The learned model nevertheless
performed substantially better under the same processed test protocol. This
difference supports the need to learn relationships among garments rather than
relying only on general visual similarity, consistent with the distinction
between similarity and complementarity made by \textcite{mcauley2015styles}.

The proposed architecture also exceeded the three controlled learned baselines
across the predefined seeds. Paired AUC intervals remained above zero for the
comparisons with the bidirectional LSTM, type-aware model and set Transformer.
Within this experiment, the result was therefore not produced by one favourable
initialisation or one unusually weak baseline run.

This does not demonstrate that the proposed model is universally superior to
the published architectures. The comparison models were controlled adaptations
using frozen SigLIP inputs and the project's four-position representation. The
evidence shows that, under these shared conditions, the lightweight design was
more effective than the selected alternatives. Direct comparison with published
headline values would not be valid because the encoders, filtering and test
coverage differ.

The proposed model was also the fastest of the four learned architectures in
the recorded training experiment. This supports its description as lightweight
within the project. The timing remains dependent on the implementation,
hardware and frozen-input pipeline and is not a universal speed comparison.

\section{Architectural evidence}

Removing the pairwise module reduced mean FITB accuracy by 10.28 percentage
points. The paired confidence interval remained entirely above zero. This
provides the clearest structural evidence in the study: explicitly comparing
garments contributed information that was not recovered by the overall outfit
summary alone.

The finding accords with the relational treatment of compatibility in earlier
work \parencite{vasileva2018typeaware,sarkar2023outfittransformer}. A garment
may be plausible in one combination but not another, so an average outfit
representation can lose differences between particular items. The pairwise
component preserves a summary of these relationships without requiring a much
larger attention model. This is evidence about the tested implementation rather
than proof that pairwise aggregation is always preferable to attention.

Removing clothing-position information produced only a small reduction, and
its confidence interval crossed zero. The experiment therefore does not provide
clear evidence that the position embedding improved FITB performance. One
possible explanation is that the visual representation and pairwise relations
already contain information that distinguishes broad garment roles. This is a
possible explanation rather than a demonstrated cause.

The position embedding was retained because it adds little computational cost
and gives the model an explicit indication of clothing role. Nevertheless, it
should not be described as an empirically necessary component. A simpler model
without it remains a reasonable alternative for future evaluation.

\section{Generalisation, overfitting and calibration}

Training AUC of 0.9699 was considerably higher than test AUC of 0.8583. The
difference shows that the model fitted the training data more strongly than it
generalised to unseen examples. The results therefore contain evidence of
training-set overfitting.

Validation and test AUC differed by only 0.00145, and the saved checkpoint
corresponded to the highest observed validation AUC. Validation-based checkpoint
selection therefore behaved consistently on the held-out test split. It does
not remove the train--test gap, but continuing beyond the selected epoch was not
supported by the validation results. This is the limited role expected of early
stopping: it controls continued fitting after validation performance peaks but
does not guarantee the absence of overfitting \parencite{prechelt1998early}.

Further tuning against the current test results would weaken the independence of
the test set. Additional regularisation experiments should therefore use
validation-only decisions or introduce a new untouched evaluation split. The
current checkpoint is a more defensible stopping point than repeatedly changing
the model after inspecting its test performance.

Temperature scaling reduced validation negative log-likelihood and expected
calibration error. The original scores were therefore more confident than
necessary on the validation data. Calibration improves the interpretation of
confidence but does not alter recommendation order or provide evidence of
improved compatibility ranking \parencite{guo2017calibration}.

\section{Outfit-length imbalance}

The dataset was balanced between positive and negative labels, so the principal
imbalance did not concern class counts. Instead, three-item outfits were more
common than four-item outfits. Moderate length weighting produced a slightly
higher overall mean AUC but reduced performance for the weakest length group and
increased the difference between the two groups.

Selecting that configuration only because its overall AUC was marginally higher
would conflict with the purpose of the weighting experiment. The predefined
gate required the weakest group to improve rather than allowing a small overall
gain to hide weaker group performance. Retaining the unweighted configuration
was therefore consistent with the stated decision rule.

This result does not show that weighting can never improve the model. It shows
only that the tested schemes did not provide a sufficiently balanced improvement
under the present data and acceptance criteria.

\section{Relationship between learned scores and explicit constraints}

The application does not allow the compatibility score to decide every part of
an outfit. This separation is necessary because statistical compatibility and
feasibility answer different questions. The model estimates whether a complete
combination resembles compatible relationships in the processed training data.
It does not know that an outer layer may be unnecessary in hot weather, that a
user-selected garment must remain fixed, or that a candidate belongs to the
requested clothing range. These conditions are applied before or during search
rather than learned indirectly from the score.

The arrangement has two analytical advantages. First, a recommendation can be
traced through candidate construction, constraint filtering and final scoring.
If no outer layer is shown, the system can identify a weather mask rather than
attribute the decision vaguely to the neural model. If an unsuitable clothing
range appears, the error lies in filtering or metadata rather than in the
compatibility score. This is a limited form of procedural explanation: it does
not explain the visual features learned by SigLIP, but it identifies which
system component made a deterministic decision.

Second, the rules stop the model from being evaluated for knowledge that its
training objective did not provide. A positive or negative outfit label does
not directly encode local weather, personal wardrobe ownership or application
state. Expecting the compatibility head to infer these factors would mix
separate sources of evidence and make failures harder to diagnose. The modular
design therefore supports auditability even though it may be less flexible
than one end-to-end model.

This division also creates limitations. A rule can exclude a combination that
the model would score highly, and errors in garment metadata can remove valid
candidates before ranking. The reported model AUC does not test those rule
failures. For this reason, model evaluation and 36-case candidate-availability
testing were reported separately. The deterministic fallback is handled in the
same way. It keeps the application usable when model artefacts are absent, but
its output should not be combined with learned-model results or described as
equivalent in quality.

\section{Candidate-search completeness and latency}

Quota-limited search examined an average of 3.41\% of the feasible combinations.
It recovered the same highest-scoring result as exact search in 27.78\% of 18
queries and recovered an average of 24.44\% of the exact top ten. The quota
stage therefore frequently removed combinations that the compatibility model
would otherwise have ranked highly.

A positive score loss in 13 of the 18 queries confirms that the change was not
limited to rearranging equally scored candidates. Score loss refers only to the
deployed model objective, however; it does not establish that a person would
prefer the exact-search result.

Exact search increased median runtime from approximately 0.028 seconds to 0.474
seconds on the recorded CPU experiment. Its 95th-percentile runtime remained
close to two seconds. The experiment identifies a clear trade-off: exact search
provides completeness under the model objective, while quota search provides
lower latency. This reflects the broader retrieval trade-off between search
coverage and computational cost \parencite{johnson2021similarity}.

Within the tested candidate sizes, the additional runtime remained compatible
with an interactive prototype. This statement applies only to the recorded
hardware and bounded search space. A substantially larger catalogue could
require approximate retrieval, staged ranking or another search strategy.

\section{Coverage, domain shift and bias limitations}

Only 29.50\% of the official FITB questions could be represented by the four
supported clothing positions. The resulting FITB accuracy is valid only for the
retained subset and cannot be presented as performance on the complete official
benchmark.

Menswear evidence was particularly sparse. Only 46 retained FITB questions
were labelled as menswear, and the training data contained very few purely
menswear outfits. The figures identify a coverage problem but do not estimate a
dependable performance difference between clothing ranges.

The feature-domain classifier separated Polyvore and Zara items with an AUC of
1.000. Their frozen feature distributions were therefore distinguishable under
the diagnostic procedure. This does not measure how much compatibility accuracy
would decline on Zara because compatible Zara outfit labels were unavailable.
The result instead prevents an unsupported assumption that performance transfers
unchanged between the two sources.

The candidate system avoids presenting historical product brands and allows the
user to choose a clothing range, retain the uploaded garment and reject
alternatives. These choices improve control and reduce some immediate product
risks, but they do not remove bias inherited from Polyvore labels or the frozen
visual representation \parencite{mehrabi2021bias,selbst2019fairness}.

Finally, no participant evidence is available. The project can support claims
about offline compatibility discrimination, candidate completion, search
completeness and system behaviour. It cannot establish emotional benefit, trust,
usability across populations or preference improvement.

\section{Summary}

The strongest model evidence comes from the controlled architecture comparison
and removal of the pairwise module. Together, these results support a lightweight
relational head on frozen visual representations within the processed Polyvore
task. Evidence for the clothing-position embedding is inconclusive.

The selected checkpoint retains an overfitting tendency, although
validation-based selection produced closely aligned validation and test results.
Calibration improved validation confidence measures, while the tested
length-weighting schemes did not provide a sufficiently balanced improvement.

Exact search substantially improved completeness under the model score at the
cost of increased runtime. The evidence remains bounded by restricted FITB
coverage, sparse menswear data, measurable feature-domain separation and the
absence of participant evaluation. Chapter 6 uses these positive and negative
findings to answer the research questions and assess the contribution.
